\section{Experimental design}
\subsection{Systems and operating points}
The dense reference is AudioLDM-M-Full~\cite{liu2023audioldm}, with $415.96$\,M U-Net parameters. We study the two released pruning severities from Singh et al.~\cite{singh2026pruning}. Severity~1 retains $145.67$\,M parameters after $65.0\%$ pruning, while severity~2 retains $71.08$\,M after $82.9\%$ pruning. At each severity, P denotes the checkpoint before recovery and \PFT{} the released checkpoint after $10^{6}$ AudioCaps fine-tuning steps. Their paired difference measures what the released recovery stage adds within the pruned architecture. It does not by itself identify an effect caused specifically by pruning.

The primary duration comparison uses $3.84$ and $10.24$\,s. Severity~2 adds $5.12$ and $7.68$\,s, producing a four-point sweep. AudioCaps \textsc{test} supplies the in-domain prompts~\cite{kim2019audiocaps}. Clotho provides a held-out sound-event domain with a different caption distribution~\cite{drossos2020clotho}. A second held-out battery uses hip-hop and rap captions from MusicCaps~\cite{agostinelli2023musiclm}. The main severity~2 AudioCaps analysis uses $192$ prompts, Clotho uses $96$, and the extended hip-hop battery uses $127$.

Generation uses DDIM~\cite{song2021ddim} with $50$ steps, guidance $2.5$ and $\eta=0$. A pre-specified robustness check repeats the severity~2 endpoint comparison with the published DDIM $200$, guidance $3.5$ recipe.

\subsection{Interventions and follow-up checks}
The mechanism and boundary checks were specified after the original operating-point result but frozen before their own generation and scoring. They therefore test explanations of the primary finding without being folded into its original confirmatory family.

The central intervention is symmetric. Starting from the same severity~2 P checkpoint, we run two full-U-Net fine-tunes for $20{,}000$ steps with the same optimizer and data recipe. One uses random $3.84$\,s AudioCaps crops and the other uses $10.24$\,s examples. Both checkpoints are evaluated at both durations on the same $192$ prompts with common generation noise. Let $J_{3.84}$ and $J_{10.24}$ denote their long-minus-short recovery interactions. Training-duration specialization predicts
\begin{equation}
\Delta J = J_{10.24}-J_{3.84}>0,
\end{equation}
because moving the training duration to $10.24$\,s should move recovery toward that evaluation duration.

A reduced dense $2\times2$ repeats the two $20{,}000$-step fine-tunes from AudioLDM-M-Full. This experiment is diagnostic rather than a substitute for the unavailable dense checkpoint from the released $10^{6}$-step recovery. The dense baseline is EMA-averaged, whereas the short experimental exports are raw weights, and the dense model had already received substantial AudioCaps training. We therefore use this $2\times2$ to test duration specialization at the same short budget, not to estimate the million-step dense recovery that no longer exists.

Additional checks extend the released recovery to $15.36$\,s on $96$ matched prompts, enlarge severity~1 with a disjoint $96$-prompt sample, and complete dense hip-hop anchors for all $127$ prompts. These checks address response shape, cross-severity stability and metric-floor interpretations.

\subsection{Pairing, endpoints and estimands}
Within each operating point, compared checkpoints receive identical diffusion noise $x_T$ for a given prompt. The prompt is the statistical unit, and $95\%$ intervals use a prompt-level percentile bootstrap with $B=10^{4}$ resamples. This common-random-number design makes $R$ a paired effect rather than a difference between independently sampled generation sets.

The primary endpoint is fused CLAP cosine from \texttt{laion/\allowbreak clap-htsat-fused}, revision \texttt{365dea6e}~\cite{wu2023clap}. Human-CLAP provides a second alignment scorer. KL is measured against matched real references, with positive recovery meaning that \PFT{} reduces KL relative to P. For PANNs, \emph{top-10 capture} is the number of ground-truth AudioSet labels from the reference clip that appear among the generated clip's ten highest-scoring PANNs classes. PANNs capture recovery is the paired increase in that count from P to \PFT{}.

For operating point $o$, recovery gain is
\begin{equation}
R(o)=\mathbb{E}_{i}\left[S(\PFT,i,o)-S(P,i,o)\right].
\end{equation}
The duration interaction is $J=R(10.24\,\mathrm{s})-R(3.84\,\mathrm{s})$. We also report
\begin{equation}
\rho_{\mathrm{ref}}=\frac{R}{S(\mathrm{ref})-S(P)},
\end{equation}
which expresses the gain as a fraction of the available gap to a dense or real-audio reference.

% ========================== FIGURE 1 SPECIFICATION ==========================
% IMPORTANT: The following figure is intentionally a layout placeholder. Do not
% replace the reserved boxes with a synthetic graph. Generate the final figure
% from the committed per-prompt result artifacts.
%
% Intended final figure footprint:
%   - one figure* spanning both columns
%   - two side-by-side panels
%   - each plot area exactly about 0.49\textwidth wide and 4.80 cm high
%   - caption below, matching the current placeholder footprint
%
% PANEL (a), operating-point response of the released recovery
%   Scientific purpose:
%     Show the main phenomenon and the uncertainty beyond the native duration.
%   X axis:
%     requested duration in seconds: 3.84, 5.12, 7.68, 10.24, 15.36.
%   Y axis:
%     paired CLAP recovery gain R = CLAP(P+FT)-CLAP(P).
%   Main series:
%     severity-2 released 10^6-step recovery at 3.84/5.12/7.68/10.24,
%     n=192, point estimate with 95% prompt-bootstrap CI.
%   Matched extension:
%     plot the n=96 matched-subset point at 10.24 s separately (open marker)
%     and connect THAT point, not the n=192 point, to the 15.36 s n=96 result
%     with a dashed segment. This visually respects the paired D4 comparison.
%   Annotation:
%     show D4 = R(15.36)-R(10.24 matched) = +0.021 [-0.023,+0.067].
%     Use wording "no clear increase" or just show the CI. Do not label plateau.
%   Optional reference cue:
%     a light vertical guide at 10.24 s may mark the native/released recovery
%     duration, but must not imply it is a peak.
%
% PANEL (b), symmetric training-duration intervention
%   Scientific purpose:
%     Visualize the falsifiable specialization prediction with matched budget.
%   X axis:
%     evaluation duration: 3.84 and 10.24 s.
%   Y axis:
%     20k-step fine-tuning gain over the SAME severity-2 P baseline.
%   Series 1:
%     shortft, trained at 3.84 s, R at both evaluation durations with 95% CIs.
%   Series 2:
%     longft, trained at 10.24 s, R at both evaluation durations with 95% CIs.
%   Do NOT include the public dense text-FT line in the main panel. It weakens
%   the symmetry and is only a supporting non-matched reference in the prose.
%   Annotate or bracket the two slopes:
%     J_short = +0.065 [+0.044,+0.087]
%     J_long  = +0.031 [+0.010,+0.052]
%     Delta J = J_long-J_short = -0.035 [-0.064,-0.004].
%   Visual logic:
%     specialization predicts the long-trained line should have a larger
%     long-minus-short slope. The observed ordering is the opposite.
%   Styling:
%     simple point-lines with vertical 95% CI whiskers, no decorative shading,
%     direct labels or a small legend, readable at two-column ICASSP print size.
% ===========================================================================
\begin{figure*}[t]
\centering
\begin{minipage}[t]{0.49\textwidth}
\centering
\fbox{\rule{0pt}{4.80cm}\rule{0.965\linewidth}{0pt}}
\end{minipage}\hfill
\begin{minipage}[t]{0.49\textwidth}
\centering
\fbox{\rule{0pt}{4.80cm}\rule{0.965\linewidth}{0pt}}
\end{minipage}
\caption{Duration dependence and the symmetric training-duration test. Left, the released severity~2 recovery is evaluated across duration, with the $15.36$\,s extension compared on its matched $96$-prompt subset. Right, two $20{,}000$-step fine-tunes from the same pruned checkpoint differ only in training duration and are evaluated at both $3.84$ and $10.24$\,s.}
\label{fig:duration}
\end{figure*}



